\documentclass[11pt,a4paper]{article}

\usepackage[utf8]{inputenc}
\usepackage[T1]{fontenc}
\usepackage{amsmath,amssymb,amsthm}
\usepackage{booktabs}
\usepackage{graphicx}
\usepackage{longtable}
\usepackage{hyperref}
\usepackage[margin=2.5cm]{geometry}
\usepackage{xcolor}

\hypersetup{colorlinks=true,linkcolor=blue!50!black,citecolor=blue!50!black,urlcolor=blue!50!black}

\newcommand{\Kj}{\ensuremath{K_j}}
\newcommand{\pp}{\,pp}
\newtheorem{proposition}{Proposition}

\title{A Certified Primal--Dual Solver for the Selective Routing\\
Problem with Synchronisation, and a Falsification Study\\
of Ten Strengthening Techniques}

\author{David Chouchena \and Tomer Tunitsky}
\date{%
  Tel-Aviv University\\[2pt]
  \emph{Exact Algorithms for Combinatorial Optimization Problems}\\[8pt]
  August 2026
}

\begin{document}
\maketitle

\begin{abstract}
The Selective Routing Problem with Synchronisation (SRPS) asks which of a set of
profitable jobs to serve, how to route parallel processors through them, and when
to start each job, subject to the requirement that every processor a job needs
must begin serving it at the same instant. We describe a certified primal--dual
solver for SRPS pairing an Adaptive Large Neighbourhood Search with a Lagrangian
relaxation of the joint-service coupling, so that every returned solution carries
a valid per-instance upper bound and hence a certificate of proximity to
optimality. On the 660-instance benchmark it attains a mean certified gap of
0.133\%, a median of 0.046\%, and proves optimality on 42.9\% of instances.

The bulk of this report is a systematic attempt to strengthen that solver using
ten techniques from the exact-algorithms syllabus, and an account of why nine
fail. Two were the primary brief: importing the cut separation logic of the exact
branch-price-and-cut method, and coupling the search more tightly to the
Lagrangian dual. Both are falsified on structural rather than parametric grounds.
Separation against the incumbent can never find a violated cut, because the
incumbent is feasible by construction; redirected at the relaxed solution, where
violations demonstrably exist, the natural feasibility cut is \emph{unsound} on
the majority of instances tested. The dual-guided operators are degenerate under
their own initialisation---the fair profit split forces the reduced profit to be
identically zero for every job---and once repaired, an isolating control
attributes the apparent gain entirely to a warm-start dual solve, which a
full-budget re-run then shows to be a screening artifact.

A ceiling argument computed from the solver's own gap distribution bounds
\emph{any} technique aimed at the certified gap at $0.1136\pp$ on average and
$0.0087\pp$ at the median---below the run-to-run variance of the experiments built
to detect it---and every gap-directed technique tested was absorbed by it. The two
aimed at runtime were not so bounded, and one delivered: parallel evaluation of
the per-processor subproblems is $1.42\times$ faster with bit-identical bounds.
Beyond the negative results we extract two methodological lessons: mechanism
\emph{activation} is not evidence of mechanism \emph{effect}, and a verdict
measured at one computational budget can reverse at another---in both directions.
\end{abstract}

\tableofcontents

\section{Introduction}\label{sec:intro}

The Selective Routing Problem with Synchronisation (SRPS) arises from configuring
the EMIR spectrograph at the Gran Telescopio Canarias~\cite{RieraLedesma2021SRPS}.
The instrument has 55 pairs of sliding bars in parallel bands; observing a target
requires moving one or more bar pairs into position, and a target lying between
bands requires several bars positioned \emph{simultaneously}. Observing time is
scarce, so only a subset of requested targets can be served. The problem is
therefore selective (which targets), routing (in what order each bar pair visits
them), and scheduling (when each observation starts), with a synchronisation
constraint binding the three together.

This report has two parts. The first describes a certified primal--dual solver:
an Adaptive Large Neighbourhood Search supplies incumbents, a Lagrangian
relaxation supplies valid upper bounds, and every solution is returned with a
certificate of how far from optimal it can possibly be. The second, and longer,
asks whether that solver can be strengthened using techniques from the
exact-algorithms syllabus.

Two directions motivated the study:

\begin{description}
  \item[Cut separation.] The exact method for SRPS derives part of its
    strength from cutting planes separated over the synchronisation structure. If
    those inequalities are valid, they should be exploitable inside a heuristic
    --- to filter infeasible candidate moves, or to tighten the dual bound.
  \item[Dual coupling.] The Lagrangian multipliers $\mu_{j,k}$ carry
    information about which jobs the relaxation considers under- or
    over-penalised. Steering the destroy and repair operators by the reduced
    profit $b_j-\sum_k \mu_{j,k}$ should focus the search where the dual says
    value remains.
\end{description}

Both are falsified. Exhausting the surrounding design space added eight further
techniques --- primal-to-dual seeding, dual stabilisation, Lagrangian
decomposition, warm-start effort scheduling, reduced-cost variable fixing,
parallel subproblem evaluation, and a standalone exact reference --- of which one
improves what it targets. Sections~\ref{sec:dual}--\ref{sec:runtime} are organised
by the metric each technique aims at, because that division is the organising fact
of the study: everything aimed at the certified gap was bounded in advance by a
ceiling computable from the solver's own results, and the two aimed at runtime
were not.

\paragraph{Contributions.}
\begin{enumerate}
  \item A certified primal--dual solver for SRPS (Section~\ref{sec:solver}) and
    its evaluation on the 660-instance benchmark (Section~\ref{sec:baseline}).
  \item A ceiling argument bounding \emph{any} gap-directed strengthening, derived
    from the solver's own gap distribution (Section~\ref{sec:ceiling}).
  \item A structural impossibility result for cut separation against a primal
    incumbent, and a soundness refutation of the natural feasibility cut when
    separation is redirected at the relaxed solution (Section~\ref{sec:cuts}).
  \item A degeneracy proof for the dual-guided operator design, a corrected
    re-test, and the isolating control that overturns its apparent success
    (Section~\ref{sec:guided}).
  \item A proof that Lagrangian decomposition of the natural split cannot improve
    the bound, closing that cell analytically rather than leaving it untested
    (Section~\ref{sec:ld}).
  \item Two methodological guards --- gate on effect rather than activation, and
    validate screen finalists at full budget --- that between them would have
    detected every failure here in minutes rather than hours
    (Section~\ref{sec:lessons}).
\end{enumerate}

\section{The Selective Routing Problem with Synchronisation}\label{sec:problem}

\subsection{Binding constraints}\label{sec:constraints}

An instance provides a job set $J=\{1,\dots,n\}$, where job $j$ yields profit
$b_j$ and requires a subset $\Kj\subseteq K$ of the processor set $K$; a
transition-time matrix $(t_{i\ell})$ covering travel and processing; and a global
horizon $L$. For each processor $k$, let $J_k\subseteq J$ denote the jobs
requiring it. A solution selects a subset of jobs, routes each processor through
its assigned jobs, and assigns start times. Three constraints bind:

\begin{enumerate}
  \item \textbf{Horizon.} Elapsed time from the start depot to the end depot may
    not exceed $L$ on any processor.
  \item \textbf{Synchronisation.} All processors in $\Kj$ must start job $j$ at
    the same instant $s_j$. For every arc $(i\to \ell)$ traversed by any
    processor, $s_\ell-s_i\ge t_{i\ell}$. Processors may idle before a
    synchronised start, and \emph{that idle time counts toward $L$}.
  \item \textbf{Consistency.} If any processor serves $j$, every processor in
    $\Kj$ must serve $j$. Partial service is forbidden.
\end{enumerate}

The objective is to maximise $\sum_{j\in J} b_j y_j$, where $y_j$ indicates
selection.

Constraints~2 and~3 are distinct requirements and are easily conflated:
consistency says \emph{which} processors must serve a job, synchronisation says
\emph{when} they must do so. A solution can satisfy either without the other.

Constraint~2 is a difference-constraint system: earliest start times are longest
paths in the arc-weighted graph induced by the routes, and a cyclic graph means no
feasible schedule exists. This makes feasibility checkable in linear time and is
the check the solver uses (Section~\ref{sec:representation}). The idle-time clause
separates ``route length $\le L$'' from ``makespan $\le L$'': routes that each fit
comfortably inside the horizon may still be jointly infeasible once synchronised
waiting is charged. Section~\ref{sec:s0} shows the Lagrangian subproblem to be
structurally blind to exactly that distinction.

\subsection{Compact formulation}\label{sec:model}

We adopt the compact arc-flow structure, denoted SRPS-1, of Riera-Ledesma and
Salazar-Gonz\'alez~\cite{RieraLedesma2021SRPS}. Let $V_k=J_k\cup\{0,n+1\}$ be
processor $k$'s node set with start depot $0$ and end depot $n+1$; let
$z_{i\ell}^k\in\{0,1\}$ indicate that processor $k$ traverses arc $(i,\ell)$,
$x_{j,k}\in\{0,1\}$ that processor $k$ serves job $j$, $y_j\in\{0,1\}$ that job
$j$ is selected, and $s_{j,k}$ the service-start time of processor $k$ at job $j$.
The objective is $\max \sum_{j\in J} b_j y_j$ subject to route structure,
\[
\sum_{\ell\in V_k\setminus\{0\}} z_{0\ell}^k = 1,\qquad
\sum_{i\in V_k\setminus\{n+1\}} z_{i,n+1}^k = 1,
\]
\[
\sum_{i\in V_k\setminus\{j\}} z_{ij}^k = x_{j,k},\qquad
\sum_{\ell\in V_k\setminus\{j\}} z_{j\ell}^k = x_{j,k}
\quad \forall j\in J_k,
\]
joint-service coupling (constraint~3),
\[
y_j \le x_{j,k} \qquad \forall j\in J,\; k\in \Kj,
\]
and timing (constraints~1 and~2),
\[
s_{\ell,k} \ge s_{i,k} + t_{i\ell} - M\bigl(1-z_{i\ell}^k\bigr),
\qquad
s_{j,k}=s_{j,k'} \;\;\forall k,k'\in \Kj,
\qquad
s_{n+1,k}-s_{0,k} \le L .
\]
The arc-time propagation doubles as subtour elimination: start times increase
along any used arc with positive transition time, so cycles are infeasible.
Section~\ref{sec:exact} solves this formulation directly as an uncoupled
reference.

\subsection{Benchmark}\label{sec:benchmark}

We use the benchmark introduced with the problem~\cite{RieraLedesma2021SRPS}. It
comprises eight families along two axes. The \emph{distance} axis distinguishes
horizontal transition costs, $|u_i-u_\ell|$, which mimic the real instrument, from
Euclidean costs on the plane; families with an \texttt{E} prefix are Euclidean.
The \emph{class} axis records $|\Kj|$: class-1 families have $|\Kj|=1$ for every
job and reduce to the team orienteering problem, while classes 2--4 have
$|\Kj|\in\{2,3,4\}$ and carry genuine synchronisation. Profits take values in
$\{1,10,100\}$. Instances are labelled by family, size $n$, and a
budget-tightness parameter $\alpha\in\{0.25,0.50,0.75\}$ controlling how much of
the total service demand fits inside $L$.

Our primary study set is the 660 instances of classes 2--4, where synchronisation
is active. Where a wider set of 780 appears (Section~\ref{sec:ceiling}), it is the
660 plus the class-1 families, and the text says so.

\section{Related work}\label{sec:related}

\subsection{SRPS and its exact treatment}

SRPS was introduced by Riera-Ledesma and
Salazar-Gonz\'alez~\cite{RieraLedesma2021SRPS} together with the model SRPS-1 of
Section~\ref{sec:model}, a benchmark, and an exact branch-price-and-cut algorithm
in two variants: one enhancing the pricing step, one node-based. The method runs
under a one-CPU-hour limit and reports a primal--dual gap of
$100\times(\text{dual}-\text{primal})/\text{primal}$.

Its reported performance defines the space in which a heuristic can add value. In
the hardest class-4 families all instances up to $n\le 55$ are solved to
optimality; at $n=60$ nonzero gaps remain, reaching 1.64\% on the Euclidean family
at $\alpha=0.75$; for $n\ge 65$ no benchmark solutions are published. Two regimes
therefore exist --- hard instances with residual exact gaps, and larger instances
with no reported exact result.

The exact method's cut separation is the direct motivation for the first of
these (Section~\ref{sec:cuts}).
Branch-price-and-cut separates violated inequalities against the \emph{fractional}
solution of a linear master problem; whether that machinery transfers to a
heuristic which never forms such a solution is precisely what
Section~\ref{sec:cuts} answers.

\subsection{Metaheuristics for selective routing}

Adaptive Large Neighbourhood Search and related destroy--repair metaheuristics are
effective across orienteering, team orienteering, and vehicle-routing variants
with rich side constraints, because they combine strong primal search with
problem-specific insertion
logic~\cite{RopkePisinger2006ALNS,Vansteenwegen2011Orienteering}. A substantial
body of metaheuristics has been developed for the team orienteering problem
specifically~\cite{Archetti2007MetaheuristicsTOP}. These deliver strong incumbents
but generally provide \emph{no per-instance optimality certificate}, which is the
gap the solver of Section~\ref{sec:solver} closes.

\subsection{Lagrangian relaxation and decomposition}

Lagrangian relaxation with subgradient optimisation is a classical route to
instance-specific bounds: a complicating coupling is dualised, the residual
problem decomposes, and the dual value bounds the
optimum~\cite{Fisher1981Lagrangian,Held1974Subgradient,Guignard2003Lagrangean}.
The approach has a long record in
routing~\cite{Fisher1997VRPTW,Kohl1999LagrangeVRPTW}, and matheuristics
interleaving such relaxations with metaheuristic search are an established way to
obtain both strong primal solutions and
bounds~\cite{ArchettiSperanza2014Matheuristic}. Step sizes follow
Polyak~\cite{PolyakStepSize}. Guignard's treatment of Lagrangian
\emph{decomposition}~\cite{Guignard2003Lagrangean} is used directly in
Section~\ref{sec:ld}, where the condition for decomposition to improve on
relaxation --- that the dualised block lack the integrality property --- decides one
cell of the design space without an experiment.

\subsection{Techniques this study draws on}

Beyond the two hypotheses, the study evaluates cutting planes and their separation
discipline, Lagrangian decomposition, reduced-cost variable fixing, dual
stabilisation by trust region and smoothing, parallel evaluation of independent
subproblems, and direct solution of the compact model by a general-purpose MIP
solver. Table~\ref{tab:coverage} maps each to the lecture that introduces it.

\section{Summary of techniques evaluated}\label{sec:coverage}

Table~\ref{tab:coverage} lists every technique evaluated, the metric it targets, and its outcome. \textbf{The numbered rows are the ten strengthening techniques} and are the subject of Sections~\ref{sec:dual}--\ref{sec:runtime}. The two unnumbered rows above the rule are not interventions: one is the solver's own dual bound, the object the ten techniques try to strengthen, and the other is an independent reference solver used for comparison in Section~\ref{sec:exact}.

The organising observation is the split by target. Interventions aimed at the certified gap are bounded in advance by the ceiling of Table~\ref{tab:ceiling}, and every one of them was absorbed by it; the two aimed at runtime carry no such bound, and one of them delivered.

\begin{table}[htbp]
\centering
\small
\caption{Techniques evaluated in this study. Rows 1--10 are the strengthening attempts; the two rows above the rule are the bound being strengthened and the reference solver. Section references point to the derivation and the supporting measurement.}
\label{tab:coverage}
\begin{tabular}{clll}
\toprule
\# & Technique & Target & Outcome \\
\midrule
--- & \emph{Lagrangian relaxation} & \emph{dual $ub$} & \emph{the method under study} (\S\ref{sec:lagbound}) \\
--- & \emph{Compact MILP (SRPS-1) in CPLEX} & \emph{reference} & \emph{looser bound than the Lagrangian on 22/30} (\S\ref{sec:exact}) \\
\midrule
1 & BPC-style cut separation & dual $ub$ & falsified on three independent grounds (\S\ref{sec:cuts}) \\
2 & Dual-guided destroy / repair & primal $z$ & $+0.0037\pp$ vs.\ warm control at production budget (\S\ref{sec:guided}, \S\ref{sec:prod}) \\
3 & Warm-$\mu$ dual effort scheduling & dual $ub$ & $-0.0086\pp$ at production budget (\S\ref{sec:prod}) \\
4 & Primal-to-dual seeding & dual $ub$ & 9 tighter / 8 looser after flooring (\S\ref{sec:seeding}) \\
5 & Dual stabilisation (trust region, smoothing) & dual $ub$ & budget-dependent; no transfer at the solver's budget (\S\ref{sec:convergence}) \\
6 & Lagrangian decomposition (variable copies) & dual $ub$ & provably equal to the existing bound (\S\ref{sec:ld}) \\
7 & Dantzig--Wolfe reformulation / column generation & dual $ub$ & master bound equals the Lagrangian dual optimum (\S\ref{sec:combining}) \\
8 & Surrogate relaxation & dual $ub$ & dominates in theory; forfeits separability (\S\ref{sec:combining}) \\
9 & Reduced-cost variable fixing & runtime & reach $24.8\%\to0.3\%$ at $0.99z$ (\S\ref{sec:fixing}) \\
10 & Parallel subproblem evaluation & runtime & $\mathbf{1.42\times}$ dual, bit-identical bounds (\S\ref{sec:parallel}) \\
\bottomrule
\end{tabular}
\end{table}


\section{The solver}\label{sec:solver}

An ALNS produces feasible solutions, a Lagrangian relaxation produces valid upper
bounds, and an outer loop alternates between them until the gap is small enough to
stop. The distinguishing feature is that the dual is not a one-off pre-processing
step but a mechanism re-tightened \emph{inside} the search, around the current
incumbent.

\subsection{Representation and feasibility}\label{sec:representation}

A solution is a selected job set and one elementary route per processor.
Feasibility is maintained by the longest-path check of
Section~\ref{sec:constraints}: the routes induce a precedence graph, earliest
start times are longest paths within it, and the solution is feasible exactly when
that graph is acyclic and its longest path from start to end depot does not exceed
$L$. This single check covers constraints 1 and 2 together, including synchronised
idle time.

\subsection{Destroy and repair operators}

Three destroy operators are used. \emph{Random destroy} removes a uniformly
sampled subset of selected jobs; \emph{worst destroy} removes the least valuable
selected jobs; \emph{Shaw destroy} removes jobs related to a random pivot under
travel-based proximity. The destroy size is sampled dynamically, with escalation
tiers enlarging the removed fraction on hard instances (Section~\ref{sec:loop}).

Three repair operators are used: \emph{profit-greedy}, \emph{ratio}, and
\emph{regret-2}. They differ only in the order in which candidate jobs are offered
to a common insertion procedure, and it is that procedure which carries the
SRPS-specific content.

\subsection{Synchronisation-aware insertion}

Inserting a job $j$ is not a single-route decision: every processor in $\Kj$ must
accommodate it, and all must start it at the same instant. The procedure
enumerates the top few local positions on each required processor route, forms
their Cartesian product, and checks each cross-route placement with the
longest-path test.

\begin{proposition}[Insertion-complexity bound]
If at most $c$ local candidate positions are retained per required processor for
job $j$, the procedure evaluates at most $c^{|\Kj|}$ route-position combinations,
each followed by one feasibility check.
\end{proposition}

\begin{proof}
Candidate positions are chosen independently on each processor in $\Kj$; their
Cartesian product contains at most $\prod_{k\in \Kj} c = c^{|\Kj|}$ combinations,
each checked once.
\end{proof}

With $c=3$ and class-4 jobs this is at most $3^4=81$ combinations per job. This
combinatorial insertion is where synchronisation enters the heuristic; the rest
follows a standard destroy--repair pattern.

\subsection{Acceptance, adaptation, and multi-start}

The search uses simulated-annealing acceptance, adaptive roulette-wheel operator
selection, and exponentially smoothed score updates: worsening moves may be
accepted to escape local optima, and operators producing improving or accepted
moves earn larger future weight. The search cycles over six fixed random seeds.

\subsection{The Lagrangian bound}\label{sec:lagbound}

The complicating constraints are the joint-service couplings $y_j \le x_{j,k}$.
Relaxing them with multipliers $\mu_{j,k}\ge 0$ gives
\begin{equation}\label{eq:lag}
L(\mu) = \sum_{j\in J} \max\Bigl(0,\; b_j - \sum_{k\in \Kj}\mu_{j,k}\Bigr)
       + \sum_{k\in K} O_k(\mu),
\end{equation}
where $O_k(\mu)$ is the optimum of a single-processor orienteering problem in
which job $j\in J_k$ carries modified profit $\mu_{j,k}$.

\begin{proposition}[Validity of the bound]
For every $\mu\ge 0$, the SRPS optimum satisfies $z^\ast \le L(\mu)$.
\end{proposition}

\begin{proof}
Any feasible SRPS solution satisfies $y_j\le x_{j,k}$ for all required processors,
so after relaxation the nonnegative penalty terms vanish on it and its objective
cannot exceed the relaxed objective at the same $\mu$. Maximising over feasible
solutions gives the claim.
\end{proof}

Each $O_k(\mu)$ is solved \emph{exactly} by bitmask dynamic programming in
$O(2^{|J_k|}|J_k|^2)$, with no pruning, sampling, or candidate-set restriction.
The exponential factor is bounded by the largest per-processor job set rather than
by $n$: since each job requires only $|\Kj|\in\{2,3,4\}$ of the $|K|=55$
processors, jobs spread thinly and each processor serves few. That the resulting
bound on $|J_k|$ is small, and does not grow with $n$, is an empirical property of
the benchmark, reported with the other measurements in
Section~\ref{sec:subproblem-size}.

Two properties of $O_k(\mu)$ matter throughout. First, it is solved to optimality,
so no looseness originates there. Second, it enforces route length $\le L$ but is
structurally blind to synchronisation idle time, because each processor is solved
in isolation. The relaxed solution is therefore optimistic in a specific,
identifiable way --- a fact Section~\ref{sec:s0} exploits and Section~\ref{sec:ld}
returns to.

\subsection{Subgradient scheme}\label{sec:subgradient}

Multipliers are initialised by the fair profit split $\mu_{j,k}=b_j/|\Kj|$, which
makes the initial relaxed objective coincide with an independent-orienteering
bound. The dual is then minimised by projected subgradient descent with
Polyak-style steps~\cite{PolyakStepSize},
\[
\eta_t = \frac{L(\mu_t)-\theta^\ast}{\|g_t\|^2},
\qquad
g_{j,k}=x^\ast_{j,k}-y^\ast_j,
\qquad
\mu_{j,k}\leftarrow \max(0,\mu_{j,k}-\eta_t g_{j,k}),
\]
clamped for stability, with a diminishing fallback schedule after 30 consecutive
non-improving iterations. The level reference $\theta^\ast$ is the current
incumbent $z$ throughout. The best bound found is returned with its multiplier
vector, retained as a warm start for the next invocation.

That the fair split is the natural \emph{starting point} for descent, yet the one
multiplier vector at which reduced profits vanish identically, is the whole of the
defect analysed in Section~\ref{sec:degeneracy}.

\subsection{The adaptive control loop}\label{sec:loop}

The destroy--repair search sits inside an outer loop allocating effort only where
certification remains hard. Each \emph{phase} runs the ALNS with cycling seeds for
up to 300 seconds. If the phase improves the incumbent, the destroy \emph{tier}
resets; otherwise the method performs an in-loop Lagrangian re-bound and escalates
the destroy fraction through 25\%, 33\%, and 40\%. Three stop conditions apply:
the upper bound matches the incumbent, the certified gap falls below 0.3\%, or the
tiers are exhausted. A per-instance cap of 3600 seconds applies throughout.

Each re-bound is capped at 200 subgradient iterations or 60 seconds, is
warm-started from the current multipliers, and uses the current incumbent as the
Polyak target. Because the numerator $L(\mu_t)-z$ shrinks as the bound approaches
the incumbent, steps are large when the bound is loose and fine when it is nearly
tight: the dual is pulled down toward the incumbent while the ALNS pushes the
incumbent up, and the two chase each other. That 200-iteration budget is the
configuration this study measures, and Section~\ref{sec:convergence} shows it to
be the binding constraint on the dual.

\subsection{Certification}\label{sec:certification}

The reported certified gap is
\begin{equation}\label{eq:gap}
\text{gap} \;=\; \frac{ub-z}{ub},\qquad ub=\lfloor L(\mu)\rfloor ,
\end{equation}
with $z$ the incumbent. Flooring is valid because all profits are integer. A gap
of $x\%$ guarantees the solution lies within $x\%$ of the optimum; a gap of $0\%$
\emph{proves} optimality by weak duality, and we call such instances \emph{proven
optimal}. No published benchmark value enters the algorithm: both the ALNS and the
dual reference use only the current incumbent.

Equation~\eqref{eq:gap} has \textbf{two levers}, and the distinction organises the
study: a mechanism can improve the gap by \emph{raising} $z$ (primal) or by
\emph{lowering} $ub$ (dual). Cutting planes act on the dual. Neighbourhood
operators act on the primal. Sections~\ref{sec:dual} and~\ref{sec:primal} divide
on exactly that line.

\section{Baseline performance}\label{sec:baseline}

\subsection{Certified gaps}

All runs use a single machine, six worker processes, and the configuration of
Section~\ref{sec:loop}. Across the 660-instance primary study set the mean
certified gap is 0.133\%, the median 0.046\%, and the maximum 1.965\%; 94.1\% of
instances are certified within 0.5\% and 100\% within 2\%. Median runtime is 5.1
minutes, and no instance stops because of the absolute runtime cap.

\subsection{Certified optima and subproblem cost}\label{sec:subproblem-size}

283 instances (42.9\%) are returned with a certified gap of exactly zero and are
therefore proven optimal in the sense of Section~\ref{sec:certification}. Of
those, 149 are also proved optimal by the exact method, where the two certificates
coincide and cross-validate one another; the remaining 134 are proven optimal by
the Lagrangian certificate alone.

The exactness of the dual bound rests on the per-processor job sets staying small,
and they do. Across all 660 instances the largest $|J_k|$ is $13$, with
per-instance means of $2.8$--$6.0$, and this maximum is \emph{flat in $n$}: it is
already $13$ at $n=80$ and remains $13$ at $n=150$. Each subproblem therefore
costs at most $2^{13}\cdot 13^2\approx 1.4\times10^{6}$ operations, so the bound is
computed exactly even on the largest instances. Section~\ref{sec:parallel} shows
the flip side of the same distribution: because cost scales as $2^{|J_k|}$, a
single processor can account for more than half an instance's dual work.

\subsection{Relation to the exact method}

Against the published branch-price-and-cut
results~\cite{RieraLedesma2021SRPS}, the instances divide into three groups. On
the 245 the exact method proves optimal, the solver recovers the optimum on 241
and is one profit unit below on four, with a mean certified gap of 0.067\%. On the
112 where the exact method returns an incumbent but stops at its one-hour limit
with a nonzero gap, the solver matches the published incumbent on 93, improves it
on 15, and is one unit below on four. On the 300 with no reported exact solution,
the solver certifies all 300, including 196 first-ever certified results.

\subsection{Headroom: a ceiling on any gap-directed improvement}\label{sec:ceiling}

The results above imply a bound on what \emph{any} technique aimed at the
certified gap can achieve, and that bound governs the rest of this report.

Cuts and every other dual-side device act only on $ub$. The optimum $P^\ast$
satisfies $z\le P^\ast\le ub$, since $z$ is attained by a feasible solution. The
maximum conceivable improvement from a technique driving the dual bound exactly to
the optimum on every instance is therefore bounded by the current certified gap
itself. The same bound applies symmetrically to primal-side techniques, which can
raise $z$ at most to $P^\ast$.

\begin{table}[htbp]
\centering
\caption{Ceiling on any gap-directed improvement, computed over 780 instances.
The row set is the 660 primary instances plus the class-1 families, which reduce
to team orienteering and are easier --- hence a proven-optimal share of 48.5\% here
against 42.9\% on the primary set. The maximum, $1.9650\pp$, is identical under
both row sets.}
\label{tab:ceiling}
\begin{tabular}{lr}
\toprule
Quantity & Value \\
\midrule
Instances already provably tight ($\text{gap}=0$) & 378 / 780 \;(48.5\%) \\
Mean maximum possible gap reduction & 0.1136\pp \\
Median maximum possible gap reduction & 0.0087\pp \\
90th percentile & 0.2948\pp \\
Maximum & 1.9650\pp \\
\midrule
Total absolute headroom $\sum(ub-z)$ & 2255 profit units \\
Share held by the top 50 instances & 45.3\% \\
\bottomrule
\end{tabular}
\end{table}

On nearly half the benchmark there is nothing to gain. On the median instance the
entire theoretical prize is $0.0087\pp$. For comparison, arm-to-arm differences in
the calibration suites below range over $\pm0.01$--$0.09\pp$: the median instance's
complete theoretical prize is smaller than the noise floor of the apparatus built
to detect it. No amount of calibration recovers a signal beneath its own
measurement noise.

This computation costs seconds and uses only data already in hand. That it was
performed \emph{after} rather than before the calibration campaigns is itself one
of the lessons of Section~\ref{sec:lessons}.

\section{Study design}\label{sec:design}

Every technique below was implemented \emph{outside} the solver: no file in the
solver's own modules was modified, and each instrument composes them from the
outside. This keeps the comparison honest --- a difference between an arm and its
control cannot be an artifact of a changed baseline.

\subsection{Notation}

\begin{table}[htbp]
\centering
\caption{Notation used throughout.}
\label{tab:notation}
\begin{tabular}{ll}
\toprule
Symbol & Meaning \\
\midrule
$z$ & incumbent objective (primal) \\
$ub$ & Lagrangian upper bound, $\lfloor L(\mu)\rfloor$ (dual) \\
$P^\ast$ & true optimum, known only to lie in $[z,ub]$ \\
$\mu_{j,k}$ & multiplier on the coupling of job $j$ to processor $k$ \\
$r_j$ & reduced profit $b_j-\sum_{k\in \Kj}\mu_{j,k}$ \\
$g_{j,k}$ & subgradient component $x_{j,k}-y_j$ \\
$G$ & absolute gap $L(\mu)-z$ \\
pp & percentage points of certified gap \\
phase & one ALNS run of up to 300\,s with cycling seeds \\
tier & destroy-fraction escalation level (25\%, 33\%, 40\%) \\
\bottomrule
\end{tabular}
\end{table}

\subsection{Instance subsets}

\begin{itemize}
  \item \textbf{Stratified-30.} A stratified subset across families B, C, D, EB,
    EC, ED and across $n$ and $\alpha$. Used wherever breadth matters.
  \item \textbf{Coupling-12.} A fast-screen subset for arm matrices comparing ten
    or more configurations.
  \item \textbf{Hard-tail-4.} The loosest-gap instances. Used where headroom is
    required for an effect to be visible at all.
  \item \textbf{Small-nonzero-30.} Every instance with $n\le 50$ carrying a
    nonzero certified gap, for the exact reference of Section~\ref{sec:exact}.
\end{itemize}

\subsection{Budget regimes}

Three regimes appear, and the distinction is load-bearing rather than incidental:
Section~\ref{sec:lessons} reports two verdicts that reverse between them.

\begin{table}[htbp]
\centering
\caption{Budget regimes. \emph{Production} is the solver's own configuration
(Section~\ref{sec:loop}); the others are screens.}
\label{tab:budgets}
\begin{tabular}{lrrrr}
\toprule
Regime & phase (s) & cap (s) & dual iters & dual (s) \\
\midrule
Production & 300 & 1800 & 200 & 60 \\
Fast screen & 120 & 240 & 120 & 20 \\
Activation probe & 20 & 45 & --- & --- \\
\bottomrule
\end{tabular}
\end{table}

The 1800\,s cap is the per-instance budget of the development harness; the solver
itself caps at 3600\,s, and no arm reached either for reasons other than tier
exhaustion.

\subsection{How each technique was tested and calibrated}

\paragraph{Cut separation (Section~\ref{sec:cuts}).}
An explicit cut pool with three separation routines (circuit, route-length,
synchronisation-path), integrated through a hard feasibility filter and a soft
scoring penalty. Calibrated by a 29-run staged suite on Stratified-30 at the
production regime: a 3-run correctness gate, 12 coarse runs spanning
low/medium/high settings, 8 refinement runs around the coarse finalists, and 6
robustness runs over seed bundles. Eight parameters were swept ---
\texttt{sep\_freq\_phase} $\in\{1,2,3\}$,
\texttt{max\_cuts\_per\_round} $\in\{5,10,20\}$,
\texttt{epsilon\_violation} $\in\{10^{-6},10^{-5},10^{-4}\}$,
\texttt{epsilon\_reject} $\in\{10^{-5},10^{-4},10^{-3}\}$,
\texttt{lambda\_cut} $\in\{0.1,0.25,0.5,1.0\}$,
\texttt{cut\_age\_limit} $\in\{2,4,8\}$,
\texttt{cut\_weight\_decay} $\in\{0,0.05,0.1\}$,
\texttt{sep\_budget\_s\_per\_phase} $\in\{5,10,20\}$ --- with lexicographic
selection on mean $\Delta$gap, then maximum final gap, then runtime.

\paragraph{Activation sweep.}
Four points of \texttt{epsilon\_violation} $\in\{10^{-7},-50,-100,-200\}$ on a
three-instance subset at the activation-probe regime, with a gate failing a run
when no cut enters the pool. Negative tolerances are non-physical and were used
deliberately to force activation, separating ``the mechanism does not apply'' from
``the mechanism is broken''.

\paragraph{Separation against the relaxed solution.}
A read-only instrument runs the subgradient to convergence (200 iterations or
30\,s), replays the relaxation at $\mu^\ast$, and measures two violation channels.
A second searches for any global ordering achieving makespan $\le L$ by time-boxed
random-restart hill climbing (90\,s per instance), reporting \textsc{schedulable}
constructively when one is found. Neither has parameters to calibrate: they are
measurements, not treatments.

\paragraph{Dual-guided coupling (Section~\ref{sec:guided}).}
Six arms on Coupling-12 at the fast-screen regime: plain control, warm-$\mu$
control, original (v1) operators with warm $\mu$, and corrected (v2) destroy,
repair, and both. The v2 repair scores $(b_j+w\,r_j)/(1+\text{cost})$ with $w=1$;
the v2 destroy scales its tie-breaking jitter to 1\% of the observed spread of
$r_j$. The decisive arm is the warm-$\mu$ control, isolating the cost of the seed
from the effect of the guidance. Three arms were repeated at the production
regime.

\paragraph{Reduced-cost fixing (Section~\ref{sec:fixing}).}
No search runs. The subgradient is run per instance on Stratified-30, $r_j$ read
off at $\mu^\ast$, and the two fixing tests evaluated in closed form. Because the
threshold $G$ depends on incumbent quality and fixing runs mid-search rather than
post hoc, reach is reported at $z$, $0.99z$ and $0.95z$.

\paragraph{Seeding and stabilisation (Sections~\ref{sec:seeding},
\ref{sec:convergence}).}
Ten variants over three axes --- initialisation
($\varepsilon\in\{0.05,0.10,0.25\}$), step rule (Polyak; trust region with radius
$0.5$ and $1.0$), and dual smoothing ($0.3$, $0.6$) --- plus two combinations, on
Stratified-30 at 200 iterations across six workers. Because a smoke test and the
full sweep disagreed, a convergence probe then ran four variants to 1000
iterations on Hard-tail-4, retaining the full bound history so the running best
could be read at checkpoints $60/100/200/400/700/1000$ from one run per variant.

\paragraph{Parallel subproblem evaluation (Section~\ref{sec:parallel}).}
The per-processor loop is farmed to a worker pool with the instance loaded once
per worker. Processors are partitioned by longest-processing-time assignment
weighted by $2^{|J_k|}$; contiguous chunking was measured first and found badly
imbalanced. Correctness requires the bound to be \emph{identical}, not merely
close, to the sequential path. Timing uses best-of-three repeats.

\paragraph{Uncoupled exact reference (Section~\ref{sec:exact}).}
SRPS-1 is built directly and handed to CPLEX 22.1.1 with a 180\,s limit and six
threads, sharing no state with the solver. Only instances CPLEX proves closed
admit a decomposition of the certified gap.

\section{Techniques targeting the dual bound}\label{sec:dual}

Everything here aims to lower $ub$, and is therefore bounded by
Table~\ref{tab:ceiling} before any of it is run.

\subsection{BPC-style cut separation}\label{sec:cuts}

\subsubsection{First design and its structural impossibility}\label{sec:cuts-structural}

The initial implementation maintained a pool of cuts of three types --- circuit,
route-length, and synchronisation-path --- separated at phase boundaries against
the ALNS incumbent, applied through a hard feasibility filter and a soft scoring
penalty.

This design cannot work, for a reason requiring no experiment. Separation is
performed against the incumbent, and the incumbent is feasible by construction.
Feasibility means precisely: no cycles in the constraint graph, no route over $L$,
and a synchronisation path within $L$ (Section~\ref{sec:constraints}). These are
exactly the three conditions the separator tests. The candidate set is therefore
empty at every call, identically, on every instance.

A 29-run calibration suite was nevertheless executed on Stratified-30, consuming
127.9 minutes and tuning eight parameters. Its outcome was a tautology: the
recommended configuration was the arm with all cut machinery disabled, and the
mean objective delta was identical ($-3.300$) across the top five arms. Every run
reported \texttt{cuts\_added\_total}~$=0$.

\subsubsection{The activation sweep and mechanism inertness}\label{sec:cuts-inert}

To separate ``the mechanism is broken'' from ``the mechanism does not apply'', a
four-point sweep drove the violation tolerance $\varepsilon$ negative, forcing
cuts into the pool.

\begin{table}[htbp]
\centering
\caption{Activation sweep on a three-instance subset. Pool size responds to
$\varepsilon$; nothing else does.}
\label{tab:sweep}
\begin{tabular}{lrrrrl}
\toprule
$\varepsilon$ & cuts added & mean viol.\ after & filter rejections & mean $\Delta$gap (pp) & final objectives \\
\midrule
$10^{-7}$ & 0  & 0.0 & 0 & $+0.057444$ & 4971 / 2285 / 2226 \\
$-50$     & 2  & 0.0 & 0 & $+0.057444$ & 4971 / 2285 / 2226 \\
$-100$    & 8  & 0.0 & 0 & $+0.057444$ & 4971 / 2285 / 2226 \\
$-200$    & 15 & 0.0 & 0 & $+0.057444$ & 4971 / 2285 / 2226 \\
\bottomrule
\end{tabular}
\end{table}

The pool grows fifteenfold; every outcome column is bit-identical to six decimal
places. That is not a statistical null --- a real treatment with no effect produces
scatter --- but an algebraic identity, and it has a specific cause. The evaluation
routine scores each cut against a right-hand side frozen at construction as the
true horizon $L$:
\[
\text{viol} \;=\; \max\bigl(0,\; \texttt{max\_len}-c.\texttt{rhs}\bigr),
\qquad c.\texttt{rhs}=L .
\]
For any feasible candidate $\texttt{max\_len}\le L$, so the clamp returns exactly
$0.0$ regardless of $\varepsilon$. The guard \texttt{if viol > 0.0} never fires, no
cut is ever recorded as violated, the hard filter never rejects, and the soft
penalty subtracts $\lambda\cdot 0$. Fifteen active cuts produced
\texttt{mean\_violation\_after}~$=0.0$.

We call this \emph{mechanism inertness}: the machinery activates and reports
activation while remaining incapable of influencing the outcome variable. It is
more dangerous than an obvious crash, because it produces clean-looking null
results that invite the interpretation ``the idea does not help'' when the correct
interpretation is ``the idea was never tested''.

A telemetry defect compounded it. A counter intended to record penalty
applications was incremented once per job removal regardless of any active cut,
reporting nonzero values on runs with an empty pool. An analyst reading only the
summary columns would have seen a mechanism apparently doing work.

\subsubsection{Redirecting separation at the relaxed solution}\label{sec:s0}

The structural argument also indicates the repair. In branch-price-and-cut,
separation runs against a \emph{fractional} solution, which is where violations
live. The heuristic has no such solution, but it does have the Lagrangian relaxed
solution, and that object is infeasible in a well-defined way: $g_{j,k}$ is
nonzero until convergence, and each $O_k(\mu)$ ignores synchronisation idle time
(Section~\ref{sec:lagbound}).

We built a read-only instrument (S0) that runs the subgradient to convergence,
replays the relaxation at $\mu^\ast$, and measures two channels:

\begin{description}
  \item[Channel A (consistency).] Jobs partially claimed --- served by some but not
    all processors in $\Kj$, violating constraint~3 directly.
  \item[Channel B (synchronisation time).] The joint difference-constraint system
    over the relaxed routes: whether it is cyclic, and if not, whether its longest
    path exceeds $L$.
\end{description}

Both fire abundantly. Notably, the same cycle-detection routine that returned zero
on every one of the 29 calibration runs returns nonzero here \emph{without
modification}. The cut family was never wrong; the object it was pointed at was.

Table~\ref{tab:s0} summarises S0 over the 30-instance stratified subset. Both violation channels fire on essentially the whole set.

\begin{table}[htbp]
\centering
\caption{S0 violation instrument over 30 instances. Every per-processor route respects the horizon, yet the joint structure does not.}
\label{tab:s0}
\begin{tabular}{lr}
\toprule
Quantity & Value \\
\midrule
Channel A: partially claimed jobs & 184 \;(184 on multi-processor jobs) \\
Mean subgradient norm $\|g\|^2$ at $\mu^\ast$ & 72.5 \\
Channel B: instances cyclic under per-processor optimal order & 21 / 30 \\
Total circuits detected & 148 \\
Acyclic instances with makespan $>L$ & 4 / 9 \\
\quad mean / max excess over $L$ & 1401.0 \;/\; 2044.0 \\
Every per-processor route within $L$ & yes \\
\bottomrule
\end{tabular}
\end{table}

The contrast in the last two rows is the substantive finding. The decomposition certifies every individual route as fitting inside the horizon, while the joint difference-constraint system does not admit a schedule under those same routes. The per-machine subproblem cannot represent the constraint that binds, because that constraint is exactly the one the relaxation dualised.


\subsubsection{Why the natural cut is unsound}\label{sec:s0b}

Channel B appears to be the opportunity: structural infeasibility the
decomposition cannot see. It does not survive scrutiny.

The relaxation fixes \emph{which} jobs each processor takes, not the \emph{order}.
Any single global total order induces per-processor sequences whose arcs all point
forward, hence an acyclic graph. Acyclicity is therefore always achievable, and
the cycles S0 reports are an artifact of each processor adopting its own locally
optimal min-time order. The real question is whether some ordering also satisfies
makespan $\le L$.

This is a soundness requirement, not a technicality. A no-good cut asserting
``this relaxed selection is infeasible'' is valid only if \emph{no} ordering works.
If a feasible ordering exists and the cut is applied anyway, a feasible point is
excluded, $ub$ becomes invalid, and every certified gap computed downstream is
wrong. For a method whose contribution is certification, that failure mode is
worse than adding no cuts at all.

Instrument S0b therefore searches for any global ordering achieving makespan
$\le L$ by time-boxed random-restart hill climbing, reporting \textsc{schedulable}
when one is found.

S0b searches for any global ordering achieving makespan $\le L$. A \textsc{schedulable} verdict is constructive: an explicit feasible ordering is exhibited, so the selection provably cannot be cut.

\begin{table}[htbp]
\centering
\caption{S0b schedulability over 30 instances. A schedulable selection cannot be excluded by a no-good cut without invalidating the bound.}
\label{tab:s0b}
\begin{tabular}{lr}
\toprule
Outcome & Instances \\
\midrule
\textsc{schedulable} (feasible ordering found) & 16 / 30 \;(53.3\%) \\
None found (inconclusive, not a proof) & 14 / 30 \;(46.7\%) \\
\bottomrule
\end{tabular}
\end{table}

On 16 of 30 instances the relaxed selection is provably schedulable, so a cut asserting its infeasibility would exclude a feasible point and invalidate every downstream certificate. The margins are narrow: the tightest feasible schedule clears the horizon by 1 time unit, and 12 of the 16 schedulable instances clear it by less than 50.

That narrowness compounds the problem. A separation routine would have to distinguish ``feasible by a handful of time units'' from ``infeasible'', and in a certifying method a misclassification does not produce a weak cut: it produces an invalid upper bound and a silently wrong certified gap. The standard discipline for user cuts is explicit on this point---a valid inequality must strengthen the relaxation \emph{without excluding feasible integer points}~\cite{Crowder1983ZeroOne,Barnhart1998BranchPrice}---and this cut family cannot meet it.


\subsubsection{Verdict}\label{sec:cuts-verdict}

We falsify cut separation on three independent grounds. First, separation
against the
incumbent cannot produce a violated cut, confirmed operationally by a 29-run suite
reporting no cuts in any run. Second, forcing activation admitted up to fifteen
cuts while leaving mean post-evaluation violation at $0.0$ and every final
objective bit-identical: the cuts were \emph{inert} rather than ineffective, and an
experiment with no power cannot support a claim about efficacy. Third, redirecting
separation at the relaxed solution yields a cut family that is \emph{unsound} on
16 of 30 instances. Independently of design, Table~\ref{tab:ceiling} caps any cut
family at $0.1136\pp$ mean and $0.0087\pp$ median.

Three limits should be stated. We do not claim BPC cuts are ineffective for SRPS
in general: we tested three families inside one architecture. The 14 of 30
instances on which no feasible ordering was found remain \emph{formally open},
since the schedulability search is heuristic; the soundness argument rests only on
the 16 constructive cases, which suffices, because a single wrongly excluded
feasible point invalidates the bound. And of the three limbs, only the ceiling
generalises beyond the tested subsets.

\subsection{Can the incumbent itself generate cuts?}\label{sec:incumbent-cuts}

A natural question, worth answering explicitly because it motivates much of the
design above: the solver always holds a best known solution, so is there
information in it that can be used to generate cuts?

The answer is no, and the reason is Section~\ref{sec:cuts-structural} stated in
general form. A cut is useful only if some point violates it. The incumbent is
feasible, so it violates no valid inequality --- that is what feasibility means.
Any separation routine run against it returns an empty candidate set, whatever the
cut family, independently of the particular families we implemented.

What the incumbent \emph{can} supply is a bound, not a cut. It is a lower bound on
the optimum, and the solver already exploits that: $z$ is the Polyak level
reference (Section~\ref{sec:subgradient}), the threshold in the fixing tests of
Section~\ref{sec:fixing}, and one end of the interval producing the ceiling of
Table~\ref{tab:ceiling}. The distinction is between a feasible point, which
bounds, and an infeasible one, which cuts. Generating cuts requires an infeasible
object, and in this architecture the only one available is the Lagrangian relaxed
solution --- which is why Section~\ref{sec:s0} redirects separation there, and
Section~\ref{sec:s0b} shows why even that fails.

\subsection{Primal-to-dual seeding and dual stabilisation}\label{sec:seeding}

The solver seeds $\mu$ by the fair split and chains warm starts from one re-bound
to the next; the incumbent enters only as the Polyak target, never as structure
(Section~\ref{sec:subgradient}). Seeding $\mu$ from the incumbent's own selection
is a genuine primal-to-dual coupling and was untested. Dual stabilisation --- trust
region and smoothing --- was likewise untested against the plain Polyak scheme.

Two coupling directions from the syllabus remained untested after the dual-guided operators: primal-to-dual seeding, in which the incumbent's own selection sets the starting multipliers rather than only serving as the Polyak target; and dual stabilisation via trust region or smoothing~\cite{DuMerle1999Stabilized}. Both are evaluated below over 30 instances at 200 subgradient iterations.

\begin{table}[htbp]
\centering
\caption{Dual variants over 30 instances, 200 iterations. $\Delta$UB is the change in the real-valued bound against baseline; negative is tighter. The final column is what survives $\lfloor\cdot\rfloor$, which is what the certified gap actually uses.}
\label{tab:variants}
\begin{tabular}{lrrr}
\toprule
Variant & mean $\Delta$UB & wins / 30 & floor: tighter / looser \\
\midrule
baseline (fair split, Polyak) & +0.000 & 0 & --- \\
incumbent seed, $\varepsilon=0.05$ & -0.415 & 14 & 6 / 8 \\
incumbent seed, $\varepsilon=0.10$ & -0.811 & 15 & 9 / 8 \\
incumbent seed, $\varepsilon=0.25$ & -0.858 & 12 & 6 / 8 \\
trust region, radius 0.5 & +12.478 & 14 & 7 / 14 \\
trust region, radius 1.0 & +1.576 & 18 & 9 / 10 \\
dual smoothing, 0.3 & +17.328 & 5 & 1 / 23 \\
dual smoothing, 0.6 & +53.899 & 0 & 0 / 29 \\
seed 0.10 + trust 0.5 & +12.616 & 13 & 6 / 14 \\
seed 0.10 + smoothing 0.3 & +13.855 & 8 & 4 / 19 \\
\bottomrule
\end{tabular}
\end{table}

Seeding produces a small mean tightening ($-0.81$ for $\varepsilon=0.10$, winning on 15 instances of 30), but the median instance moves by $-0.000$---that is, not at all. The mean is carried by 10 instances with $|\Delta \text{UB}|\ge 1$, spanning $-10.0$ to $+3.0$. After flooring, the direction is a coin flip, and two of the three seeding variants are net worse. There is no usable improvement here.

Smoothing is straightforwardly harmful and monotone in its weight. The trust region requires more care, and is treated separately below, because its verdict depends on the iteration budget.


\subsection{Does the trust region survive a longer budget?}\label{sec:convergence}

At 60 iterations the trust region appeared decisively ahead ($-8.07$ mean). At 200 it was behind ($+1.58$). That sign reversal is the signature of a starved baseline, the same failure this study documents for the screening budget, so the question had to be settled at a budget where the baseline is allowed to finish. Each variant was therefore run once at 1000 iterations with its bound history retained, and the running best read off at checkpoints.

\begin{table}[htbp]
\centering
\caption{Bound trajectories to 1000 subgradient iterations. \texttt{prod} is the bound the production configuration reports, which uses at most 200 iterations warm-started.}
\label{tab:convergence}
\begin{tabular}{llrrrrrr}
\toprule
Instance & Variant & 60 & 100 & 200 & 400 & 700 & 1000 & prod \\
\midrule
\texttt{ED\_n050\_011\_a50\_032} & baseline & 1308.00 & 1307.45 & 1306.70 & 1306.30 & 1306.30 & 1306.30 & 1307 \\
\texttt{} & trust05 & 1319.10 & 1311.08 & 1306.15 & 1306.01 & 1306.00 & 1306.00 & 1307 \\
\texttt{} & trust10 & 1309.37 & 1307.12 & 1306.28 & 1306.07 & 1306.05 & 1306.05 & 1307 \\
\texttt{} & seed\_eps10 & 1308.00 & 1307.81 & 1307.16 & 1306.31 & 1306.31 & 1306.31 & 1307 \\
\midrule
\texttt{EB\_n110\_006\_a50\_017} & baseline & 3334.37 & 3315.45 & 3283.32 & 3257.70 & 3254.17 & 3253.54 & 3259 \\
\texttt{} & trust05 & 3339.00 & 3332.44 & 3320.78 & 3300.94 & 3275.49 & 3254.66 & 3259 \\
\texttt{} & trust10 & 3335.55 & 3327.42 & 3305.58 & 3271.27 & 3253.84 & 3253.10 & 3259 \\
\texttt{} & seed\_eps10 & 3347.50 & 3327.50 & 3273.31 & 3256.03 & 3254.12 & 3253.46 & 3259 \\
\midrule
\texttt{B\_n140\_021\_a25\_061} & baseline & 3314.00 & 3272.00 & 3234.77 & 3199.73 & 3194.34 & 3192.64 & 3195 \\
\texttt{} & trust05 & 3365.35 & 3335.34 & 3289.67 & 3228.13 & 3204.08 & 3190.54 & 3195 \\
\texttt{} & trust10 & 3333.36 & 3295.92 & 3233.06 & 3203.97 & 3188.54 & 3187.46 & 3195 \\
\texttt{} & seed\_eps10 & 3313.70 & 3270.75 & 3230.83 & 3197.30 & 3193.42 & 3191.90 & 3195 \\
\midrule
\texttt{D\_n080\_041\_a50\_122} & baseline & 2243.01 & 2242.34 & 2240.95 & 2238.51 & 2236.60 & 2235.94 & 2235 \\
\texttt{} & trust05 & 2238.84 & 2237.82 & 2236.42 & 2235.11 & 2234.58 & 2234.46 & 2235 \\
\texttt{} & trust10 & 2241.62 & 2239.68 & 2237.71 & 2236.26 & 2235.15 & 2234.81 & 2235 \\
\texttt{} & seed\_eps10 & 2243.11 & 2242.07 & 2240.21 & 2237.32 & 2235.97 & 2235.40 & 2235 \\
\bottomrule
\end{tabular}
\end{table}

At 1000 iterations the trust region leads on all four instances, so the 200-iteration verdict was premature rather than the 60-iteration one being spurious: the variants had simply not separated yet. After flooring, 2 of the four show a strictly tighter integer bound (B\_n140\_021\_a25\_061 by 5, D\_n080\_041\_a50\_122 by 1).

Two effects must be separated before claiming anything, the same decomposition that exposed the warm-$\mu$ confound. Part of the gain against production is simply running the subgradient longer: baseline alone improves several instances once given 1000 iterations rather than 200. Only the residual is attributable to stabilisation. On \texttt{B\_n140} the bound moves from 3195 (production) to 3192 (baseline at 1000) to 3187 (trust region at 1000), so of eight integer units, five come from the budget and three from the trust region.

The practical conclusion is budget-conditional and therefore does not transfer to the solver as configured. It caps the dual at 200 iterations or 60 seconds, warm-started and invoked only when a phase fails to improve---and at 200 iterations the trust region is worse on average and 9-to-10 after flooring. Its advantage exists only at roughly five times that budget, on the hardest instances, at a proportionate runtime cost. We therefore report it as a scoped observation about the dual in isolation, not as a recommended change.


\subsection{Alternative ways to combine the subproblem solutions}\label{sec:combining}

Equation~\eqref{eq:lag} combines the per-processor optima by plain summation:
$L(\mu)$ is the $y$-term plus $\sum_k O_k(\mu)$. That is the only aggregation the
solver performs, and it is worth asking what alternatives exist and what they
would buy. Three are available, and all three can be settled without an
experiment.

\paragraph{Lagrangian decomposition.} Duplicate the shared variables and price
their disagreement rather than summing independent blocks. For the natural split
this provably reproduces the same bound, because the dualised block has the
integrality property; Section~\ref{sec:ld} gives the argument. A different split,
separating routing from synchronisation timing, is not covered by it and remains
the one open direction.

\paragraph{Column generation and Dantzig--Wolfe.} Reformulate with one column per
per-processor route and price columns by the same orienteering subproblem. By the
Dantzig--Wolfe equivalence, the linear master's bound at convergence equals the
Lagrangian dual optimum $\min_\mu L(\mu)$ for the same subproblem structure. Column
generation therefore cannot yield a \emph{tighter} bound than a fully converged
subgradient; it can only reach it differently. That is not nothing --- an LP master
does not oscillate, and Section~\ref{sec:convergence} shows the subgradient still
descending when its budget expires --- but the improvement on offer is one of
convergence, not of bound quality, and it is bounded by the same ceiling.

\paragraph{Surrogate relaxation.} Aggregate the coupling constraints into a single
weighted surrogate rather than dualising each. The surrogate bound dominates the
Lagrangian bound in theory, but the surrogate subproblem retains the aggregated
constraint and is no longer separable by processor --- forfeiting precisely the
property that makes $O_k(\mu)$ exactly solvable. The aggregation that would
tighten the bound is the one that destroys the decomposition.

The pattern across all three matches the ceiling result: the aggregation is not
where the looseness lives. The looseness lives in the subproblem's blindness to
synchronisation waiting (Section~\ref{sec:lagbound}), and no way of combining
blind subproblem solutions restores what the blindness discarded.

\subsection{Warm-start dual effort scheduling}\label{sec:warm}

The isolating control introduced for the dual-guided operators
(Section~\ref{sec:guided}) doubles as a test of a dual-side technique in its own
right: spending a subgradient solve \emph{before} the search rather than only
reactively when a phase stalls. Its result is reported with those arms, because
the two are measured together, and Section~\ref{sec:prod} shows it to be a
screening artifact.

\section{Techniques targeting the primal incumbent}\label{sec:primal}

\subsection{Dual-guided destroy and repair}\label{sec:guided}

\subsubsection{Design and the initialisation degeneracy}\label{sec:degeneracy}

Two operators were added to the ALNS portfolio. A guided destroy removes jobs
ranked by ascending reduced profit $r_j$, breaking ties on descending dual
pressure. A guided repair inserts jobs by a score derived from $r_j$ and the
insertion cost. Multipliers are optionally refreshed from the in-loop re-bound.

Multipliers are initialised by the fair profit split
(Section~\ref{sec:subgradient}). Then
\[
\sum_{k\in \Kj}\mu_{j,k}\;=\;|\Kj|\cdot\frac{b_j}{|\Kj|}\;=\;b_j
\qquad\Longrightarrow\qquad
r_j \;=\; b_j-\sum_{k\in \Kj}\mu_{j,k} \;=\; 0
\]
for \emph{every} job, identically. Direct measurement confirms it: on one 80-job
instance, $0$ of $80$ jobs have nonzero reduced profit; on a 150-job instance, $0$
of $150$.

The fair split is the correct starting point for descent, but it is precisely the
multiplier vector at which the guidance signal vanishes. Two consequences follow
mechanically:

\begin{description}
  \item[Destroy degenerates to random removal.] The sort key is
    $(0+\text{uniform}(-10^{-6},10^{-6}),\,-\text{dual pressure})$. The first
    component is pure noise on an identically zero base, and exact ties have
    measure zero, so the intended second criterion never breaks anything.
  \item[Repair degenerates to insertion order.] With $r_j=0$ every job takes a
    constant fallback score, the strict improvement test fails for every candidate
    after the first, and insertion follows the iteration order of the unserved
    set.
\end{description}

Both degenerate operators are inserted at position~0 of their portfolios,
displacing tuned operators.

\subsubsection{The original results are what the defect predicts}

\begin{table}[htbp]
\centering
\caption{Original dual-guided matrix (30 instances). Lower mean $\Delta$gap is
better.}
\label{tab:h2-orig}
\begin{tabular}{lrrl}
\toprule
Arm & mean $\Delta$gap (pp) & mean $\Delta$obj & Predicted by \\
\midrule
baseline     & 0.000313 & $-0.033$ & --- \\
full         & 0.003921 & $-0.067$ & both defects, partly rescued by feedback \\
destroy only & 0.008313 & $-0.200$ & random destroy displacing a tuned operator \\
repair only  & 0.016217 & $-0.367$ & insertion-order repair displacing regret-2 \\
\bottomrule
\end{tabular}
\end{table}

The ordering is not noise: it is what the degeneracy analysis predicts in advance.
The repair-only arm, displacing regret-2 with order-dependent insertion, is worst;
the destroy-only arm, displacing a tuned operator with randomness, is next; the
baseline wins.

\subsubsection{Correcting the design}\label{sec:h2-fix}

Testing the hypothesis rather than the defect requires two repairs. \textbf{Warm
multiplier initialisation} seeds $\mu$ from a real subgradient solve: nonzero
reduced profits rise from $0/80$ to $74/80$, with spread $9.667$.
\textbf{Non-degenerate operators (v2)} let the dual signal \emph{modulate} a
competent base heuristic instead of replacing it, so that with $r\equiv 0$ the
repair reduces to profit-density greedy rather than to a constant. Both original
operators are retained behind flags, so the degenerate behaviour remains
reproducible.

\subsubsection{Corrected results}\label{sec:h2-results}

Table~\ref{tab:h2fixed} reports the corrected matrix, in which the dual signal is live at the point of use (Section~\ref{sec:h2-fix}).

\begin{table}[htbp]
\centering
\caption{Corrected dual-guided matrix. Lower mean $\Delta$gap is better. The $\delta$ column is measured against the warm-$\mu$ control where available, so the seed cost is not charged to the guidance.}
\label{tab:h2fixed}
\begin{tabular}{lrrrr}
\toprule
Arm & $n$ & mean $\Delta$gap (pp) & mean $\Delta$obj & $\delta$ vs.\ base (pp) \\
\midrule
control (no dual operators, no warm $\mu$) & 12 & 0.1663 & -0.250 & +0.1310 \\
control + warm $\mu$ (isolates seed cost) & 12 & 0.0353 & -0.250 & --- \\
v1 operators + warm $\mu$ & 12 & 0.0559 & -0.750 & +0.0206 \\
v2 destroy only + warm $\mu$ & 12 & 0.0316 & -0.417 & -0.0037 \\
v2 repair only + warm $\mu$ & 12 & 0.0478 & -0.667 & +0.0125 \\
v2 destroy + repair + warm $\mu$ & 12 & 0.0324 & -0.333 & -0.0030 \\
\bottomrule
\end{tabular}
\end{table}

The matrix decomposes cleanly once the warm-$\mu$ control is included. Warm initialisation alone accounts for -0.1310\pp of the movement; the guided operators, measured against that same baseline, span -0.0037\pp to +0.0206\pp, and 2 of the 4 are worse than the warm control outright.

Two columns settle the attribution. The mean objective delta is identical for the plain and warm controls (-0.250 in both), so warm initialisation changed nothing in the primal---its entire contribution is a tighter $ub$. The guided operators, by contrast, degrade the objective (to -0.750) while moving the gap by at most 0.0037\pp in the favourable direction, inside the $\pm0.01$--$0.09\pp$ arm-to-arm noise band of the earlier suites and below the $0.0087\pp$ median ceiling of Table~\ref{tab:ceiling}.

The naive reading of this matrix---comparing guided arms against the \emph{plain} control---reports a 0.1347\pp improvement and concludes that corrected dual guidance works. It does not. That apparent effect is the warm-$\mu$ solve, and it would have been reported as a positive result had the isolating arm not been run. We regard this as the most instructive artifact of the study: a confound that survives every check except the one that varies the suspected cause on its own.

The hypothesis is therefore falsified as stated, while the experiment yields a positive finding of its own: dual effort spent \emph{before} the search, rather than reactively when a phase fails to improve, tightens the certificate materially on this subset. This is a scheduling question about the dual, not a steering question about the primal, and it belongs to the dual-stabilisation family~\cite{DuMerle1999Stabilized} rather than to guidance. Section~\ref{sec:threats} records the budget caveat that qualifies it.


\subsubsection{Production-budget validation}\label{sec:prod}

The fast screen runs at \texttt{abs\_cap}~$=240$s and \texttt{lag\_max\_time}~$=20$s, against production defaults of $1800$s and $60$s---13\% of the runtime and 33\% of the dual budget. Since the effect isolated above is a \emph{dual} effect, the screen may simply have starved the baseline. We therefore re-ran the decisive arms at full budget.

\begin{table}[htbp]
\centering
\caption{Production-budget validation on the same 12 instances. The warm-$\mu$ effect does not survive a properly-fed dual.}
\label{tab:prod}
\begin{tabular}{lrrrr}
\toprule
Arm & $n$ & mean $\Delta$gap (pp) & mean $\Delta$obj & mean final gap (\%) \\
\midrule
control (no dual operators, no warm $\mu$) & 12 & 0.0061 & -0.083 & 0.2636 \\
control + warm $\mu$ & 12 & -0.0026 & +0.083 & 0.2550 \\
v2 destroy + repair + warm $\mu$ & 12 & 0.0012 & +0.000 & 0.2587 \\
\bottomrule
\end{tabular}
\end{table}

The warm-$\mu$ effect measures -0.0086\pp at production budget, against $-0.1310\pp$ in the screen---a loss of 93\% of its magnitude. The mechanism is visible in the control itself: its mean final gap falls from $0.4239\%$ under the screen to 0.2636\% here. The screen cut instances off before the dual converged, so an extra upfront subgradient solve supplied effort the baseline was otherwise denied. Given a full budget, the existing in-loop re-bounding supplies that effort anyway, and the advantage disappears.

Guidance measured against the warm control at production budget is +0.0037\pp---still nothing, and still on the wrong side of zero.

This is the study's second false-signal mode, and the more dangerous one in practice. Reduced-budget screening is standard for parameter exploration, and here it produced a confident, plausible, and entirely spurious $0.13\pp$ result. Unlike the inert-cut null, nothing internal to the screen betrays it: the arms are consistent, the ordering is stable, and the effect is far outside run-to-run noise. Only re-running the finalist at full budget exposes it.


\section{Techniques targeting runtime}\label{sec:runtime}

Runtime carries no ceiling. Both techniques here were selected for that reason.

\subsection{Reduced-cost variable fixing}\label{sec:fixing}

Lagrangian reduced-cost fixing is the one syllabus technique whose target is runtime rather than gap, and runtime carries no ceiling. Writing $G=L(\mu)-z$, forcing a single $y_j$ shifts the bound by one term only:

\[
L(\mu \mid y_j{=}1)=L(\mu)+\min(0,r_j),\qquad
L(\mu \mid y_j{=}0)=L(\mu)-\max(0,r_j),
\]

so $r_j<-G$ proves $j$ absent from every optimal solution and $r_j>G$ proves it present. Both tests are $O(1)$ per job given $r_j$, which the dual already computes.

\begin{table}[htbp]
\centering
\caption{Reach of reduced-cost fixing over 30 instances, as a function of incumbent quality. Fixing runs mid-search, not post-hoc.}
\label{tab:fixing}
\begin{tabular}{lr}
\toprule
Incumbent used for $G$ & Mean jobs fixable \\
\midrule
final $z$ (post-hoc, near-optimal) & 24.8\% \\
$0.99\,z$ & 0.3\% \\
$0.95\,z$ & 0.0\% \\
\midrule
Overall (final $z$) & 481 / 2490 \;(19.3\%) \\
Median per instance & 0.7\% \\
Instances with any fixable job & 15 / 30 \\
\bottomrule
\end{tabular}
\end{table}

The rule has no practical reach. Backing the incumbent off by one per cent collapses the mean from 24.8\% to 0.3\%, and by five per cent to 0.0\%. The median instance yields 0.7\% even post-hoc.

The reason is instructive, and it inverts an argument that looked compelling beforehand. We expected an asymmetry: the same tightness that caps cut-based strengthening should make the fixing threshold $G$ easy to clear. It does---observed $G$ runs as low as 0.0000---but a small $G$ does not create fixing opportunity. It reports that the instance is already closed, at which point there is no remaining search to accelerate. Both properties are consequences of the same quantity, so they cannot be played against each other.


\subsection{Parallel subproblem evaluation}\label{sec:parallel}

Profiling attributes 84\% of dual wall time to \texttt{orienteering\_dp\_with\_selection}, invoked once per processor per subgradient iteration. The per-processor solves are independent by construction~\cite{Fisher1981Lagrangian} --- that independence is what the relaxation buys --- so the loop parallelises. Unlike every other intervention tested here, this one cannot alter the objective, the bound, or validity---only wall-clock---so it cannot be falsified in the same way.

\begin{table}[htbp]
\centering
\caption{Parallel dual, 4 instances, 4 workers, 200 iterations, best of 3 repeats. Timing spread across repeats averaged 1.4\%.}
\label{tab:parallel}
\begin{tabular}{lrrrr}
\toprule
Instance & sequential (s) & parallel (s) & speedup & \% of ceiling \\
\midrule
\texttt{B\_n140\_021\_a25\_061} & 14.36 & 10.13 & 1.42$\times$ & 75\% \\
\texttt{EB\_n110\_006\_a50\_017} & 6.53 & 6.07 & 1.08$\times$ & 72\% \\
\texttt{D\_n080\_041\_a50\_122} & 14.46 & 10.09 & 1.43$\times$ & 44\% \\
\texttt{ED\_n050\_011\_a50\_032} & 0.47 & 0.27 & 1.73$\times$ & 43\% \\
\midrule
\textbf{mean} & & & \textbf{1.42}$\times$ & 59\% \\
\bottomrule
\end{tabular}
\end{table}

Bounds were identical to the sequential path in every configuration and every repeat, which is the correctness claim the approach needs.

Two structural limits are worth recording. First, DP cost is $O(2^m m^2)$ in a processor's candidate count $m$, and $m$ ranges from 0 to 13 across the 55 processors, so cost is extremely skewed: on \texttt{B\_n140} a single processor carries 53\% of the instance's total dual work. Because no partition splits one DP, speedup is capped at $\text{total}/\text{max\_single}$ regardless of worker count---a mean ceiling of 2.66$\times$ here. Second, contiguous chunking is badly wrong for such a distribution (observed imbalance up to 5.35$\times$ at eight workers); longest-processing-time assignment recovers most of the difference, lifting one instance from 1.13$\times$ to 1.71$\times$.

The end-to-end value is nonetheless small, because the component is not on the critical path. In the production-budget arms the dual was invoked on one instance of twelve; the others terminated on the gap threshold before any phase failed to improve. On the single instance that did use it, the dual accounted for roughly 45 of 1297 seconds, or 3.5\% of runtime, so a 1.42$\times$ component speedup yields about 1.01$\times$ overall and exactly 1.00$\times$ elsewhere. The honest summary is a correct and measured component improvement whose component does not matter much in the current architecture; it would matter in a cold-start regime where the initial bound must be computed rather than loaded.


\section{An uncoupled exact reference}\label{sec:exact}

The ceiling argument bounds any strengthening by $(ub-z)/ub$ because the optimum $P^\ast$ is known only to lie in $[z,ub]$. Solving instances exactly would pin $P^\ast$ down and convert that bound into a measurement. We therefore solved the compact SRPS-1 arc-flow model directly in CPLEX~22.1.1---standalone, sharing no state with the ALNS or the Lagrangian---on the 30 instances with $n\le 50$ that still carry a nonzero certified gap. Instances already certified optimal are excluded: their decomposition is trivially zero.

CPLEX closed 7 of 30 within a 180-second limit. Only those admit a valid decomposition; on a timed-out run CPLEX's incumbent is not $P^\ast$, and on 10 of the 23 unproven instances it was strictly worse than the ALNS incumbent.

\begin{table}[htbp]
\centering
\caption{The 7 instances CPLEX closed. \emph{looseness} is $(ub-P^\ast)/ub$, recoverable by a tighter dual; \emph{subopt} is $(P^\ast-z)/ub$, which no dual improvement can reach.}
\label{tab:exact}
\begin{tabular}{lrrrrr}
\toprule
Instance & $P^\ast$ & $z$ & $ub$ & looseness & subopt \\
\midrule
\texttt{ED\_n040\_002\_a75\_006} & 1026 & 1026 & 1029 & 0.29\% & 0.00\% \\
\texttt{ED\_n040\_004\_a50\_011} & 1164 & 1154 & 1165 & 0.09\% & 0.86\% \\
\texttt{D\_n045\_007\_a25\_019} & 662 & 662 & 663 & 0.15\% & 0.00\% \\
\texttt{C\_n050\_004\_a50\_011} & 2017 & 2016 & 2017 & 0.00\% & 0.05\% \\
\texttt{EC\_n050\_004\_a50\_011} & 1885 & 1885 & 1886 & 0.05\% & 0.00\% \\
\texttt{C\_n050\_004\_a25\_010} & 1242 & 1242 & 1243 & 0.08\% & 0.00\% \\
\texttt{D\_n050\_012\_a25\_034} & 753 & 753 & 754 & 0.13\% & 0.00\% \\
\bottomrule
\end{tabular}
\end{table}

The ALNS is \emph{provably optimal} on 5 of these 7, so on those the entire certified gap is bound looseness and nothing is left for the primal to recover. Across the unproven instances CPLEX matched the ALNS on 12 and beat it on 1, which is external verification of incumbent quality from a solver sharing no code with the heuristic.

The more useful finding concerns the bound. On 22 of 30 instances the Lagrangian certificate is \emph{strictly tighter} than the dual bound CPLEX reaches within its limit. The theory anticipates this: the Lagrangian subproblem is an orienteering problem solved exactly by dynamic programming, and orienteering lacks the integrality property, so its Lagrangian bound may strictly dominate the LP relaxation of the compact model. It does here. This is also why the literature attacks SRPS by branch-price-and-cut rather than solving SRPS-1 directly: the compact relaxation is weak.

The decomposition the experiment was built for is nonetheless limited, for a structural reason worth recording. A direct MILP reaches only small instances, and small instances are largely closed already---75\% of those with $n\le 50$ have a certified gap of exactly zero---leaving 30 candidates of which CPLEX closed 7. The instances where the split would matter carry gaps above 1.5\% at $n=110$--$150$, far beyond reach when the specialised branch-price-and-cut of the literature closes only $n\le 60$. On the hard tail the ceiling remains an analytic bound; it cannot be replaced by measurement.


\section{Lagrangian decomposition: why the natural split cannot help}\label{sec:ld}

One cell of the design space calls for Lagrangian \emph{decomposition} rather than
the Lagrangian \emph{relaxation} the solver implements. We did not run it, because
for the natural split it can be shown not to help.

Write the problem as $\max\{c^\top v : v \in X_A \cap X_B\}$, where $X_A$ collects
the per-processor routing and horizon constraints and $X_B$ the joint-service
coupling $x_{j,k}=y_j$, both intersected with integrality. The solver dualises
$X_B$ and solves over $X_A$ exactly, so its bound is
\[
z_{\mathrm{LR}}=\max\{c^\top v : v \in \operatorname{conv}(X_A),\; x_{j,k}=y_j\}.
\]
Lagrangian decomposition duplicates the shared variables, assigns one copy to each
block, and prices their disagreement:
\[
z_{\mathrm{LD}}=\max\{c^\top v : v \in \operatorname{conv}(X_A)\cap
\operatorname{conv}(X_B)\}.
\]
By the Guignard--Kim result~\cite{Guignard2003Lagrangean},
$z_{\mathrm{LD}}\le z_{\mathrm{LR}}$ always, with strict improvement exactly when
$\operatorname{conv}(X_B)$ is strictly contained in the linear description of
$X_B$ --- that is, when the dualised block lacks the integrality property.

Here it does not lack it. $X_B$ is the set of binary vectors satisfying
$x_{j,k}=y_j$: equalities tying each $x_{j,k}$ to a single binary $y_j$. Fixing
$y$ determines $x$, and the remaining polytope is a product of $\{0,1\}$-valued
coordinates, so every vertex of the linear description has $y_j\in\{0,1\}$. Hence
$\operatorname{conv}(X_B)$ coincides with that description, and
\[
z_{\mathrm{LD}}=z_{\mathrm{LR}} .
\]
Duplicating variables for the coupling constraint is guaranteed to reproduce the
bound the solver already computes. No experiment is required, and none would be
informative: the two bounds are equal by construction, so any measured difference
could only be numerical noise.

\paragraph{The split that could help.}
A different partition is not covered by this argument. Separating \emph{routing}
from \emph{synchronisation timing} --- one block enforcing per-processor routes, the
other the difference-constraint system on start times --- dualises a copy constraint
whose blocks both have non-integral hulls, so $z_{\mathrm{LD}}<z_{\mathrm{LR}}$
becomes possible. This is also the split targeting a deficiency we measured rather
than a hypothetical one: the per-processor subproblem enforces route length but is
blind to synchronisation waiting (Sections~\ref{sec:constraints},
\ref{sec:lagbound}).

We record it as future work. It requires a second subproblem family, a timing
oracle, and a second multiplier vector, and its payoff remains capped by
Table~\ref{tab:ceiling}. It is the only remaining idea in this study with a
mechanism pointing at a measured source of looseness, which is why we name it
precisely rather than leaving the cell blank.

\section{Methodological lessons}\label{sec:lessons}

\paragraph{Activation is not effect.}
A mechanism can fire, report firing, and still be incapable of changing the
outcome. The activation sweep of Table~\ref{tab:sweep} passed an explicit
activation gate at $\varepsilon=-100$ while remaining completely inert. Gates must
therefore be placed on the \emph{outcome} variable: before spending a calibration
suite, assert on a small subset that the treated arm differs from control. That
guard would have caught the failure of Section~\ref{sec:cuts} in roughly fifteen
minutes instead of 128.

\paragraph{Budgets decide verdicts, in both directions.}
Reduced-budget screening is standard practice, and it produced here a confident,
internally consistent, entirely spurious $0.13\pp$ result. Nothing inside the
screen betrays it: arm ordering is stable and the effect sits far outside
run-to-run noise. The failure mode is specific --- a starved baseline flatters any
treatment that supplies the starved resource.

The same mechanism runs the other way, and we walked into it. The trust region
measured $-8.07$ at 60 iterations, $+1.58$ at 200, and led on every instance at
1000. We declared it falsified on the middle number. It was not falsified; the
arms had not yet separated. A negative verdict taken at a single budget is as
unsafe as a positive one. Three of this study's findings --- the fast screen, the
trust region, and reduced-cost fixing --- turn on this, which suggests it is the
dominant confound in metaheuristic calibration rather than an incidental one.

\paragraph{Check the ceiling before calibrating.}
Table~\ref{tab:ceiling} costs seconds and uses only data already in hand. It
establishes that the median instance offers $0.0087\pp$ of headroom --- below
experimental noise --- which is decisive information about whether a calibration
campaign can possibly succeed. It should precede the campaign, not follow it.

\paragraph{Verify that the signal is non-degenerate.}
The defect of Section~\ref{sec:degeneracy} is a single arithmetic identity,
$\sum_k b_j/|\Kj|=b_j$, checkable in one line before any experiment. Guidance
mechanisms should assert that their signal has nonzero variance across the
decision set \emph{at the point of use}.

\paragraph{Distinguish soundness failures from performance failures.}
The result of Section~\ref{sec:s0b} is not that the cut performs poorly, but that
it is \emph{invalid} on a majority of instances. In a certifying method, an
invalid cut does not degrade the result --- it corrupts the guarantee, silently.
Soundness must be established before efficacy is measured.

\section{Threats to validity}\label{sec:threats}

The schedulability search of Section~\ref{sec:s0b} is heuristic. A
\textsc{schedulable} verdict is constructive and therefore sound: an explicit
feasible ordering is exhibited. A ``none found'' verdict is not a proof of
infeasibility, and we do not treat it as one; those instances remain formally
open, and settling them would require an exact feasibility oracle over the
disjunctive orderings, which is NP-hard in general.

The S0 and S0b samples are subsets of the benchmark rather than its entirety,
selected to include both the stratified subset and the loosest-gap instances. The
ceiling argument carries no such limitation, being computed over 780 instances,
and it is the argument on which the principal conclusion rests.

The corrected dual-guided arms carry a runtime confound by construction: warm
initialisation spends an additional subgradient solve the plain control does not
pay. We control for it with a dedicated arm applying warm initialisation with both
guided operators disabled. A second confound, the reduced screening budget, is not
merely controlled but \emph{measured}: Section~\ref{sec:prod} re-runs the decisive
arms at full budget.

All results are measured against the solver configured as in
Section~\ref{sec:loop}, in particular a dual refresh capped at 200 iterations or
60 seconds. Section~\ref{sec:convergence} shows that budget to be the binding
constraint on the dual: the subgradient is still descending when it stops. Raising
it tightens the certificate on every hard-tail instance tested, by $0.29\pp$ on
average at essentially no runtime cost. Conclusions about gap-directed techniques
become \emph{more} secure under such a configuration, since a tighter $ub$ shrinks
the ceiling and leaves any technique less room; conclusions that are themselves
budget-conditional --- the trust region above all --- would need re-testing.

Finally, we tested three specific cut families and one dual-guidance design. We do
not claim that no cut or coupling mechanism can help SRPS heuristics in general.
We claim that these do not, that the reasons are structural, and that the ceiling
bounds what any successor could achieve on this benchmark.

\section{Conclusions}\label{sec:conclusions}

We described a certified primal--dual solver for SRPS and then tried, ten times
over, to strengthen it with techniques from the exact-algorithms syllabus. Nine
attempts failed, and the reasons are structural rather than parametric.

\textbf{Cut separation fails three times over.} Separation against a feasible
incumbent is
mathematically incapable of finding a violated cut. Redirected at the Lagrangian
relaxed solution it finds violations in abundance, but the natural feasibility cut
is unsound on the majority of instances tested, because the relaxed selection
admits a feasible synchronised schedule once ordering is chosen globally rather
than per processor. And independently of design, the certified bound is already
exact on 48.5\% of instances, capping every cut family at $0.11\pp$ mean and
$0.009\pp$ median improvement.

\textbf{Dual-guided coupling fails twice over.} The operator design is
degenerate under its own
initialisation, and correcting that defect was worth doing because it converts an
untested hypothesis into a tested one. The corrected matrix then appears positive,
and is not: an isolating control attributes the movement to the warm-start dual
solve rather than to guidance, and a full-budget re-run shows that residue to be a
screening artifact as well. Guidance measured against a properly-fed control is
$+0.0037\pp$ --- on the wrong side of zero.

Of the remaining techniques, Lagrangian decomposition of the natural split is
provably equivalent to the relaxation already in use; seeding and stabilisation
fail to move the floored bound; and reduced-cost fixing, though sound and cheap,
fires only once the incumbent is already essentially optimal, at which point no
search remains to accelerate. The expected asymmetry --- that the tightness capping
cut-based strengthening should make the fixing threshold easy to clear --- does not
hold, because both properties are consequences of the same quantity $G$ and cannot
be played against each other.

Parallelising the per-processor subproblems is the one intervention that improves
what it targets: $1.42\times$ on the dual with bit-identical bounds, capturing
59\% of a ceiling itself limited because a single subproblem can carry half an
instance's work. Its end-to-end value is nonetheless about one per cent, because
the dual is invoked on one instance in twelve. A correct optimisation of something
that is not the bottleneck is still not a contribution to the solver, and we
report it as such.

The broader contribution is diagnostic. Mechanism inertness --- machinery that
activates, reports activation, and cannot affect the outcome --- is a failure mode
producing clean, publishable-looking null results from experiments that never took
place. The guard against it is cheap and general: gate on the outcome variable,
never on whether the mechanism fired. And no verdict, positive or negative, should
be believed at a single computational budget.

\appendix

\section{Reproduction and results}\label{app:repro}

All instruments are read-only with respect to the solver and modify no existing
code path. The accompanying repository ships a single runner:

\begin{verbatim}
python reproduce.py --list          # every stage, its command, its runtime
python reproduce.py --all           # everything, about 6 hours
python reproduce.py --all --quick   # everything at reduced budgets
python reproduce.py --fragments     # rebuild this paper's tables from the CSVs
\end{verbatim}

Every table in Sections~\ref{sec:dual}--\ref{sec:exact} is generated directly from
the result files rather than transcribed. Artifacts are written to
\texttt{results/bpc\_replica\_dev/} (cut separation and its instruments),
\texttt{results/dual\_guided\_dev/} (coupling, seeding, stabilisation, fixing,
parallel), and \texttt{results/exact/} (the CPLEX reference). Result files are
timestamped and never overwritten, so repeated runs accumulate.

The per-instance baseline results --- incumbent, bound, certified gap, and runtime
for all 780 instances --- are in \texttt{results/adaptive\_master.csv}. Every
$\Delta$ reported in this paper is measured against that table, so any arm can be
recompared against the baseline directly.

The 71 instance files the experiments touch are bundled with the repository under
CC0; the full benchmark is not needed to reproduce anything reported here.
Repository documentation covers setup, the meaning of every result column, and the
integer-floor convention of Equation~\eqref{eq:gap}, which is the single most
common source of misreading: a technique can improve the real-valued bound and
change the certified gap by nothing.

\begin{thebibliography}{9}

\bibitem{RieraLedesma2021SRPS}
J.~Riera-Ledesma, J.-J. Salazar-Gonz\'alez,
Selective routing problem with synchronization,
Computers \& Operations Research 135 (2021) 105465.
\href{https://doi.org/10.1016/j.cor.2021.105465}{doi:10.1016/j.cor.2021.105465}.

\bibitem{RopkePisinger2006ALNS}
S.~Ropke, D.~Pisinger,
An adaptive large neighborhood search heuristic for the pickup and delivery
problem with time windows,
Transportation Science 40~(4) (2006) 455--472.
\href{https://doi.org/10.1287/trsc.1050.0135}{doi:10.1287/trsc.1050.0135}.

\bibitem{Vansteenwegen2011Orienteering}
P.~Vansteenwegen, W.~Souffriau, D.~Van Oudheusden,
The orienteering problem: A survey,
European Journal of Operational Research 209~(1) (2011) 1--10.
\href{https://doi.org/10.1016/j.ejor.2010.03.045}{doi:10.1016/j.ejor.2010.03.045}.

\bibitem{Archetti2007MetaheuristicsTOP}
C.~Archetti, A.~Hertz, M.~G. Speranza,
Metaheuristics for the team orienteering problem,
Journal of Heuristics 13~(1) (2007) 49--76.
\href{https://doi.org/10.1007/s10732-006-9004-0}{doi:10.1007/s10732-006-9004-0}.

\bibitem{Fisher1981Lagrangian}
M.~L. Fisher,
The Lagrangian relaxation method for solving integer programming problems,
Management Science 27~(1) (1981) 1--18.
\href{https://doi.org/10.1287/mnsc.27.1.1}{doi:10.1287/mnsc.27.1.1}.

\bibitem{Held1974Subgradient}
M.~Held, P.~Wolfe, H.~P. Crowder,
Validation of subgradient optimization,
Mathematical Programming 6~(1) (1974) 62--88.
\href{https://doi.org/10.1007/BF01580223}{doi:10.1007/BF01580223}.

\bibitem{Guignard2003Lagrangean}
M.~Guignard,
Lagrangean relaxation,
Top 11~(2) (2003) 151--200.
\href{https://doi.org/10.1007/bf02579036}{doi:10.1007/bf02579036}.

\bibitem{PolyakStepSize}
B.~T. Polyak,
Minimization of unsmooth functionals,
USSR Computational Mathematics and Mathematical Physics 9~(3) (1969) 14--29.
\href{https://doi.org/10.1016/0041-5553(69)90061-5}{doi:10.1016/0041-5553(69)90061-5}.

\bibitem{Fisher1997VRPTW}
M.~L. Fisher, K.~O. J\o rnsten, O.~B.~G. Madsen,
An optimization algorithm for the vehicle routing problem with time window
constraints,
Operations Research 45~(3) (1997) 395--406.
\href{https://doi.org/10.1287/opre.45.3.395}{doi:10.1287/opre.45.3.395}.

\bibitem{Kohl1999LagrangeVRPTW}
N.~Kohl, J.~Desrosiers, O.~B.~G. Madsen, M.~M. Solomon, F.~Soumis,
2-path cuts for the vehicle routing problem with time windows,
Transportation Science 33~(1) (1999) 101--116.
\href{https://doi.org/10.1287/trsc.33.1.101}{doi:10.1287/trsc.33.1.101}.

\bibitem{ArchettiSperanza2014Matheuristic}
C.~Archetti, M.~G. Speranza,
A survey on matheuristics for routing problems,
EURO Journal on Computational Optimization 2~(4) (2014) 223--246.
\href{https://doi.org/10.1007/s13675-014-0030-7}{doi:10.1007/s13675-014-0030-7}.

\bibitem{Barnhart1998BranchPrice}
C.~Barnhart, E.~L. Johnson, G.~L. Nemhauser, M.~W.~P. Savelsbergh, P.~H. Vance,
Branch-and-price: Column generation for solving huge integer programs,
Operations Research 46~(3) (1998) 316--329.
\href{https://doi.org/10.1287/opre.46.3.316}{doi:10.1287/opre.46.3.316}.

\bibitem{Crowder1983ZeroOne}
H.~Crowder, E.~L. Johnson, M.~W. Padberg,
Solving large-scale zero-one linear programming problems,
Operations Research 31~(5) (1983) 803--834.
\href{https://doi.org/10.1287/opre.31.5.803}{doi:10.1287/opre.31.5.803}.

\bibitem{DuMerle1999Stabilized}
O.~du Merle, D.~Villeneuve, J.~Desrosiers, P.~Hansen,
Stabilized column generation,
Discrete Mathematics 194~(1--3) (1999) 229--237.
\href{https://doi.org/10.1016/S0012-365X(98)00213-1}{doi:10.1016/S0012-365X(98)00213-1}.

\end{thebibliography}

\end{document}
